\subsection{Size-distribution treatment}
\label{sec:size}

Aerosol particle size distributions (APSDs) provide size-resolved
context for the PM mass and aerosol optical properties associated
with each collection category. We use number, projected-area, and
volume distributions to distinguish the contribution of smaller
particles to number concentration from the greater weighting of
larger particles in area and volume. These distributions describe
the aerosol particles sampled during the collection periods;
they do not identify individual particles as plant-material or dust.

Combining the two instruments requires an explicit relationship
between their diameter scales. The differential mobility particle
sizer (DMPS) reports electrical mobility diameter, which is retained
without conversion \citep{wiedensohler2012}. The aerodynamic particle
sizer (APS) reports aerodynamic diameter, \(D_a\). Following the
approach used at Hyyti\"al\"a by \citet{ylivinkka2025}, we convert
\(D_a\) to an approximate mobility-equivalent diameter, \(D_m\), using
\begin{equation}
D_m \approx D_a\sqrt{\frac{\rho_0}{\rho_p}},
\qquad
\rho_0 = 1~\mathrm{g\,cm^{-3}},
\qquad
\rho_p = 1.5~\mathrm{g\,cm^{-3}},
\end{equation}
where \(\rho_0\) is the reference density and \(\rho_p\) is the assumed
particle density. This spherical-particle approximation neglects
the diameter dependence of the Cunningham slip correction in the
conversion. Particle density and shape were not measured for each
collection, so the converted diameters are model-dependent
equivalent sizes rather than independently measured geometric
diameters.

The merged distributions use DMPS bins below
\(0.6~\mu\mathrm{m}\) and converted APS bins at and above
\(0.6~\mu\mathrm{m}\), avoiding duplicate use of the overlapping
instrument ranges. The upper selection criterion is applied to
the original APS aerodynamic bin centres, retaining only centres
at or below \(15~\mu\mathrm{m}\). The largest retained APS bin
centre is therefore \(14.9~\mu\mathrm{m}\) in aerodynamic diameter,
corresponding to approximately \(12.2~\mu\mathrm{m}\) after the
conversion above. This value locates that bin on the converted
diameter axis; it is not a sharp upper limit on the diameters of
all particles represented by the bin. Although the plotted axis
extends to \(15~\mu\mathrm{m}\), no values are extrapolated beyond
the retained bin centres. SI Section S9 gives the native
instrument grids.

The upper-bin selection limits the plotted range but does not
demonstrate complete particle detection within it.
\citet{vasilatou2023} documented size-dependent counting
efficiencies and differences among four APS instruments in
laboratory tests extending beyond \(10~\mu\mathrm{m}\).
Their findings motivate caution when interpreting the large-particle
end of APS distributions; they do not independently validate the
counting efficiency of our instrument or establish
\(15~\mu\mathrm{m}\) as a universal performance threshold.
This plotting criterion does not alter the independently
measured super-PM$_{10}$ mass.

Number, projected-area, and volume distributions are calculated
on the same retained diameter grid. We define
\begin{equation}
n(D) = \frac{\mathrm{d}N}{\mathrm{d}\log_{10}D},
\end{equation}
where \(N\) is particle number concentration and \(D\) denotes
mobility diameter for DMPS bins or the approximate
mobility-equivalent diameter for converted APS bins. Assuming
spherical particles, the corresponding distributions of projected
area, \(A_{\mathrm{proj}}\), and volume, \(V\), are
\begin{align}
\frac{\mathrm{d}A_{\mathrm{proj}}}{\mathrm{d}\log_{10}D}
&= \frac{\pi D^2}{4}\,n(D),\\
\frac{\mathrm{d}V}{\mathrm{d}\log_{10}D}
&= \frac{\pi D^3}{6}\,n(D).
\end{align}
The geometric factors are evaluated at the retained bin centres.
Projected area is the circular cross-sectional area of a sphere,
not its total surface area, which is four times larger
\citep{seinfeld2016}; neither quantity is an optical scattering
cross-section. Multiplying all APS diameters by the same constant
preserves logarithmic bin widths and hence the number-distribution
ordinates, while changing the diameter-dependent area and volume
weights.

Collection averaging preserves the sampling interval as the unit
of comparison. Hourly spectra are required to contain finite
values in all supplied size bins and are averaged over each
recorded filter interval using the duration of their overlap
with that interval as weights. A collection is retained when
complete spectra cover at least \(75\,\%\) of its duration;
this selection is independent of PM mass validity. The retained
collection-mean spectra then receive equal weight in each
category mean. Consequently, longer collections do not
automatically receive greater weight in the category curves.

The five supplied scalar metrics describe broader particle-size
regimes rather than individual bins in the plotted distributions.
They comprise Aitken-regime number concentration,
\(N_{\mathrm{Ait}}\), accumulation-regime area,
\(S_{\mathrm{acc}}\), accumulation-regime volume,
\(V_{\mathrm{acc}}\), coarse-regime volume,
\(V_{\mathrm{coarse}}\), and the coarse-to-accumulation volume
ratio, \(R_V\). These quantities are analysed as supplied
regime-level descriptors, separately from the newly constructed
size-resolved curves; they are not obtained by integrating the
category-mean curves. The two analyses therefore answer
complementary questions: the scalars compare specified regimes,
whereas the curves show where differences occur across particle
size. Neither the descriptive mean curves nor their geometric
weightings establish whole-curve statistical significance or
plant-material-specific and dust-specific contributions to PM mass.
